% SGA 3，报告 VII_B，英文重建本——第 1 轮（中文译本）。
% 范围：依据原文件第 11--14 页，从 0.3.5 至 0.3.8；止于原文件第 14 页的 0.4 之前。
% 合并阅读版页码： 524--527. 再版印刷页码： 513--516.
% 本范围内的原版页边定位： 485--488.
% 依据： Polo--Gille Exp7B-13oct24.pdf,
% SHA-256 D303690E0195654BBE7C3DAA5D89EA23E24CD709A3575217234AC19AE8836694.
% 曾参阅 Jacob Reinhold 的 CC BY 4.0 Markdown，作为对比版本来源。

\medskip
\noindent\textbf{0.3.5.}
\phantomsection
\label{sga3:VIIB:0.3.5}
考虑三个伪紧 \(A\)-模 \(L,M,N\)，其中 \(N\) 长度有限。
赋予 \(A\)-模 \(\operatorname{Hom}_c(M,N)\) 离散拓扑。每个元素
\[
  \psi\in
  \operatorname{Hom}_c
    \bigl(L,\operatorname{Hom}_c(M,N)\bigr)
\]
定义一个连续双线性映射
\[
  \psi':L\times M\longrightarrow N,
  \qquad
  (\ell,m)\longmapsto\psi(\ell)(m).
\]
由此得到自然同构
\begin{equation}
  \operatorname{Hom}_c
    \bigl(L,\operatorname{Hom}_c(M,N)\bigr)
  \xrightarrow{\sim}
  \operatorname{Hom}_c
    \bigl(L\widehat{\otimes}_A M,N\bigr).
  \tag{1}
  \label{sga3:VIIB:0.3.5-eq1}
\end{equation}
这是 \(\operatorname{Hom}_c(L\widehat{\otimes}_A M,N)\) 的另一种刻画。
现在将它用于 \(M\) 是伪紧 \(A\)-模的滤过投射系 \((M_i)\) 的投射极限的情形。
由 (1) 和 \SGARefLink{sga3:VIIB:0.2.F}{0.2.F}~\SGARefLink{sga3:vii:paragraph:5-1:equation:ii}{(ii)}，有自然同构
\begin{align}
  \operatorname{Hom}_c
    \left(L\widehat{\otimes}_A\varprojlim_iM_i,N\right)
  &\simeq
  \operatorname{Hom}_c
    \left(L,\operatorname{Hom}_c(\varprojlim_iM_i,N)\right)
  \notag\\
  &\simeq
  \operatorname{Hom}_c
    \left(L,\varinjlim_i\operatorname{Hom}_c(M_i,N)\right).
  \tag{2}
  \label{sga3:VIIB:0.3.5-eq2}
\end{align}

此外，由于 \(\varinjlim_i\operatorname{Hom}_c(M_i,N)\) 是离散的，
从 \(L\) 到这个模的每个连续映射都经过 \(L\) 的一个有限长度商分解。
因此，自然映射
\begin{equation}
  \varinjlim_i
  \operatorname{Hom}_c
    \bigl(L,\operatorname{Hom}_c(M_i,N)\bigr)
  \xrightarrow{\sim}
  \operatorname{Hom}_c
    \left(L,\varinjlim_i\operatorname{Hom}_c(M_i,N)\right)
  \tag{3}
  \label{sga3:VIIB:0.3.5-eq3}
\end{equation}
是同构。最后，再由 (1) 和 \SGARefLink{sga3:VIIB:0.2.F}{0.2.F}~\SGARefLink{sga3:vii:paragraph:5-1:equation:ii}{(ii)} 得到自然同构
\begin{align}
  \varinjlim_i
  \operatorname{Hom}_c
    \bigl(L,\operatorname{Hom}_c(M_i,N)\bigr)
  &\simeq
  \varinjlim_i
  \operatorname{Hom}_c
    \bigl(L\widehat{\otimes}_A M_i,N\bigr)
  \notag\\
  &\simeq
  \operatorname{Hom}_c
    \left(\varprojlim_i
      (L\widehat{\otimes}_A M_i),N\right).
  \tag{4}
  \label{sga3:VIIB:0.3.5-eq4}
\end{align}
结合 \SGARefLink{sga3:VIIB:0.3.5-eq2}{(2)}、\SGARefLink{sga3:VIIB:0.3.5-eq3}{(3)} 和
\SGARefLink{sga3:VIIB:0.3.5-eq4}{(4)}，得到下列结果。

\medskip
\noindent\textbf{命题。}
\phantomsection
\label{sga3:VIIB:0.3.5-proposition}
{\itshape 设 \((M_i)\) 是 \(\operatorname{PC}(A)\) 中对象的滤过投射系，
\(L\in\operatorname{PC}(A)\)，且 \(N\in\operatorname{LF}(A)\)。
则有关于 \(N\) 函子性的同构
\[
  \operatorname{Hom}_c
    \left(L\widehat{\otimes}_A\varprojlim_iM_i,N\right)
  \simeq
  \operatorname{Hom}_c
    \left(\varprojlim_i(L\widehat{\otimes}_A M_i),N\right),
\]
从而有同构
\[
  L\widehat{\otimes}_A\varprojlim_iM_i
  \simeq
  \varprojlim_i(L\widehat{\otimes}_A M_i).
\]
因此，完备化张量积与滤过投射极限交换。\footnote{这也重新得到了 \SGARefLink{sga3:VIIB:0.3.1.1}{0.3.1.1}：
\[
  L=L\widehat{\otimes}_A A
  =
  L\widehat{\otimes}_A
    \varprojlim_{\mathfrak l}(A/\mathfrak l)
  \simeq
  \varprojlim_{\mathfrak l}L/\mathfrak lL,
\]
其中 \(\mathfrak l\) 遍历 \(A\) 的开理想。}}

% 依据： Exp7B 原文件页码 12; 合并阅读版页码 525; 印刷页码 514.
\medskip
\noindent\textbf{0.3.6.}
\phantomsection
\label{sga3:VIIB:0.3.6}
特别地，\footnote{参见\ 编者注~(12)。} 完备化张量积与任意积交换。例如，由于 \(A\) 是其局部分量
\(A_{\mathfrak m}\) 的积（\SGARefLink{sga3:VIIB:0.1.1}{0.1.1}），每个伪紧 \(A\)-模
\[
  M\simeq A\widehat{\otimes}_A M
\]
等同于以下各模的积
\[
  M_{\mathfrak m}
  =
  A_{\mathfrak m}\widehat{\otimes}_A M,
\]
这些模称为 \(M\) 的局部分量。

同样，设 \(M\) 和 \(N\) 是伪紧 \(A\)-模。回顾 \SGARefLink{sga3:VIIB:0.2.2}{0.2.2}，有
\[
  M^\dagger=\operatorname{Hom}_c(M,A).
\]
考虑映射
\[
  \varphi:
  M^\dagger\otimes_A N^\dagger
  \longrightarrow
  (M\widehat{\otimes}_A N)^\dagger
\]
它将 \(f\otimes g\) 映到从 \(M\widehat{\otimes}_A N\) 到 \(A\) 的映射
\[
  m\widehat{\otimes}n\longmapsto f(m)g(n)
\]
。当 \(M\simeq A\) 时，映射 \(\varphi\) 是双射。

\medskip
\noindent\textbf{推论。}
\phantomsection
\label{sga3:VIIB:0.3.6-corollary}
\textit{若 \(A\) 是 Artin 环，则只要 \(M\) 是拓扑自由的，或更一般地是投射对象，
\(\varphi\) 就是同构。}

事实上，固定 \(N\) 时，由前述结果和 \SGARefLink{sga3:VIIB:0.2.F}{0.2.F}，两个函子
\[
  M\longmapsto(M\widehat{\otimes}_A N)^\dagger
  \qquad\text{和}\qquad
  M\longmapsto M^\dagger\otimes_A N^\dagger
\]
都将任意积变为直和。

\medskip
\noindent\textbf{注记 0.3.6.A。}\footnote{编者添加了本注记，它将在 \SGARefLink{sga3:VIIB:2.3.1}{2.3.1} 中用到。}
\phantomsection
\label{sga3:VIIB:0.3.6.A}
\textit{同样应用 \SGARefLink{sga3:VIIB:0.2.F}{0.2.F}，得到下列结果。设 \(A\) 是 Artin 环，
\(M,Q\in\operatorname{PC}(A)\)，且 \(N\in\operatorname{LF}(A)\)。
若 \(Q\) 是投射对象，则有自然同构}
\[
  \operatorname{Hom}_c(M,Q)
  \xrightarrow{\sim}
  \operatorname{Hom}_A(Q^\dagger,M^\dagger),
  \qquad
  Q^\dagger\otimes_A N
  \xrightarrow{\sim}
  \operatorname{Hom}_c(Q,N).
\]

\medskip
\noindent\textbf{0.3.7.}
\phantomsection
\label{sga3:VIIB:0.3.7}
对每个 \(\mathfrak m\in\Upsilon(A)\)，函子
\[
  M\longmapsto A_{\mathfrak m}\widehat{\otimes}_A M
\]
显然正合。由于每个投射伪紧 \(A\)-模 \(P\) 都是形如
\(A_{\mathfrak m}\) 的模的积，得到
\[
  M\longmapsto P\widehat{\otimes}_A M
\]
只要 \(P\) 是投射对象就是正合的。逆命题也成立。

\medskip
\noindent\textbf{命题。}
\phantomsection
\label{sga3:VIIB:0.3.7-proposition}
{\itshape 设 \(A\) 是伪紧环，\(P\) 是伪紧 \(A\)-模。下列条件等价：
\begin{enumerate}
  \item[\textup{(i)}]
  \(P\) 是 \(\operatorname{PC}(A)\) 的投射对象；
  \item[\textup{(ii)}]
  每个局部分量 \(P_{\mathfrak m}\) 都是拓扑自由 \(A_{\mathfrak m}\)-模；
  \item[\textup{(iii)}]
  函子 \(M\mapsto P\widehat{\otimes}_A M\) 正合。
\end{enumerate}}

\SGARefLink{sga3:vii:paragraph:5-1:equation:i}{(i)} 与
\SGARefLink{sga3:vii:paragraph:5-1:equation:ii}{(ii)} 的等价性由
\SGARefLink{sga3:VIIB:0.2.F}{0.2.F}~\SGARefLink{sga3:VIIB:0.2.F:item:iii}{(iii)} 和
\SGARefLink{sga3:VIIB:0.2.1}{0.2.1} 得到，而我们刚才已经证明了
\SGARefLink{sga3:VIIB:0.2.1:item:ii}{(ii)}~\(\Rightarrow\)~\SGARefLink{sga3:VIIB:0.2.F:item:iii}{(iii)}。
设函子 \(M\mapsto P\widehat{\otimes}_A M\) 正合。由于
\(P\widehat{\otimes}_A M\) 是其局部分量的积，
\[
  (P\widehat{\otimes}_A M)_{\mathfrak m}
  \simeq
  P_{\mathfrak m}
    \widehat{\otimes}_{A_{\mathfrak m}}M_{\mathfrak m},
\]
我们归约到 \(A\) 为局部环的情形。下面证明此时 \(P\) 是拓扑自由的。

设 \(\mathfrak m\) 是 \(A\) 的极大理想。商
\(P/\overline{\mathfrak m}P\) 是 \(A/\mathfrak m\) 上的线性紧向量空间，
从而是 \(A/\mathfrak m\) 的副本的积（\SGARefLink{sga3:VIIB:0.2.2}{0.2.2}）。
因此，存在一族 \(A\) 的副本 \((A_i)_{i\in I}\) 以及同构
\[
  \phi:
  \prod_{i\in I}(A_i/\mathfrak m)
  \xrightarrow{\sim}
  P/\overline{\mathfrak m}P.
\]
由于 \(\prod_{i\in I}A_i\) 是投射对象，存在交换方块
\begin{center}
  \phantomsection
  \label{sga3:VIIB:0.3.7-projective-lift-square}
  % SGA3-VII-DIAG-082
  % SGA3-VII-DIAG-082
% Authority: Exp7B-13oct24.pdf, local p. 13.
% Crop: figures/Exp7B_localp013_projective_lift_square.png
% Locator: tex/components/20_expose_VIIB_sections035_038_tensor_projectives_en.tex, Loop 1 line 260.
\begin{tikzcd}[row sep=3em,column sep=4.8em]
\displaystyle\prod_{i\in I}A_i \arrow[r,"\psi"] \arrow[d,"p"']
  & P \arrow[d,"q"] \\
\displaystyle\prod_{i\in I}(A_i/\mathfrak m)
  \arrow[r,"\varphi","{\sim}"']
  & P/\mathfrak mP
\end{tikzcd}

\end{center}
其中 \(p\) 和 \(q\) 是典范投影。将 Nakayama 引理应用于
\(\operatorname{Coker}\psi\)。由于
\[
  (A/\mathfrak m)\widehat{\otimes}_A\psi
\]
恰为 \(\phi\)，映射 \(\psi\) 是满射。\footnote{这证明了编者注~(22) 中预告的结果：每个伪紧 \(A\)-模 \(M\) 都是某个拓扑自由
\(A\)-模的商。由于 \(\operatorname{PC}(A)\) 中的积正合，可以归约到此处处理的局部情形。}

% 依据： Exp7B 原文件页码 13; 合并阅读版页码 526; 印刷页码 515.
令
\[
  B=\prod_{i\in I}A_i,
  \qquad
  N=\ker\psi.
\]
于是有下列交换正合图：
\begin{center}
  \phantomsection
  \label{sga3:VIIB:0.3.7-nakayama-exact-diagram}
  % SGA3-VII-DIAG-083
  % SGA3-VII-DIAG-083
% Authority: Exp7B-13oct24.pdf, local p. 13.
% Crop: figures/Exp7B_localp013_nakayama_exact_diagram.png
% Locator: tex/components/20_expose_VIIB_sections035_038_tensor_projectives_en.tex, Loop 1 line 285.
\begin{tikzcd}[row sep=2.05em,column sep=1.65em]
{} & {} & {} &
  0
    \arrow[d]
    \arrow[ddlll,dashed,rounded corners=10pt,
      to path={
        -- ([xshift=6em]\tikztostart.east)
        |- ([yshift=.8em]\tikztotarget.north)
        -- (\tikztotarget)
      }]
  & {} \\
{} & \mathfrak m\widehat\otimes_A N \arrow[r] \arrow[d]
  & \mathfrak m\widehat\otimes_A B \arrow[r] \arrow[d]
  & \mathfrak m\widehat\otimes_A P \arrow[r] \arrow[d]
  & 0 \\
0
  \arrow[r]
  \arrow[dr,dashed,out=-90,in=180,looseness=1.35]
  & A\widehat\otimes_A N \arrow[r] \arrow[d]
  & A\widehat\otimes_A B \arrow[r,"\psi"] \arrow[d]
  & A\widehat\otimes_A P \arrow[r] \arrow[d]
  & 0 \\
{} & (A/\mathfrak m)\widehat\otimes_A N \arrow[r]
  & (A/\mathfrak m)\widehat\otimes_A B \arrow[r,"\varphi"]
  & (A/\mathfrak m)\widehat\otimes_A P \arrow[r]
  & 0 \qquad .
\end{tikzcd}

\end{center}
对前两行应用蛇引理，表明最下面一行中的态射
\[
  (A/\mathfrak m)\widehat{\otimes}_A N
  \longrightarrow
  (A/\mathfrak m)\widehat{\otimes}_A B
\]
是单态射。但 \(\phi\) 是同构，所以 \((A/\mathfrak m)\widehat{\otimes}_A N=0\)。
由 \SGARefLink{sga3:VIIB:0.3.3}{0.3.3} 得到 \(N=0\)，从而 \(\psi\) 是同构。\footnote{使用“幂零 Nakayama 引理”的完全类似的证明表明，若 \(A\) 是 Artin 环，且
\(P\) 是平坦 \(A\)-模，则 \(P\) 的每个局部分量都是自由 \(A\)-模
（参见\ \textup{[BAC]}，II，\S~\SGARefLink{sga3:VIIB:3.2:prop}{3.2}，
命题~5 的\SGARefLink{sga3:VIIB:0.3.5-eq2}{推论~2}）。}

\medskip
\noindent\textbf{0.3.8. 推论。}
\phantomsection
\label{sga3:VIIB:0.3.8}
\textit{设 \(A\) 是完备 Noether 局部环，\(P\) 是伪紧 \(A\)-模。
则 \(P\) 是拓扑自由的，当且仅当 \(P\) 在 \(A\) 上平坦。}

事实上，典范映射
\[
  M\otimes_A P\longrightarrow M\widehat{\otimes}_A P
\]
在 \(M=A\) 时为双射，因而在 \(M\) 为 Noether 模时也为双射：
取 \(M\) 的一个有限呈示，并利用张量积和完备化张量积的右正合性即可。

现在，\(P\) 平坦，当且仅当
\[
  M\longmapsto M\otimes_A P
\]
在 Noether 模上正合。同样，\SGARefLink{sga3:VIIB:0.3.7}{命题
0.3.7} 的证明表明，只要序列
\[
  0\longrightarrow
  \mathfrak m\widehat{\otimes}_A P
  \longrightarrow
  A\widehat{\otimes}_A P
  \longrightarrow
  (A/\mathfrak m)\widehat{\otimes}_A P
  \longrightarrow0
\]
正合，\(P\) 就是拓扑自由的。因此，\(P\) 拓扑自由，当且仅当函子
\[
  M\longmapsto M\widehat{\otimes}_A P
\]
在 Noether 模上正合。由上面建立的等式
\[
  M\otimes_A P=M\widehat{\otimes}_A P
\]
即得推论。

% 下一断点： Exp7B 原文件页码 14, 合并阅读版页码 527,
% 再版印刷页码 516, 0.4 处。
